% ---------------------------------------------------------------
\section{Introduction}
\label{sec:introduction}
% ---------------------------------------------------------------

National research evaluation systems face an enduring challenge: how to assess the quality, influence, and impact of scholarly output across diverse institutions equitably, transparently, and without imposing unsustainable administrative burdens. In Australia, the Excellence in Research for Australia (ERA) framework, established in 2010 by the Australian Research Council \citep{arc2019era}, relied on comprehensive institutional data submissions evaluated through a hybrid of peer review panels and citation benchmarks. Over five rounds, ERA successfully mapped the national research landscape but incurred substantial direct compliance costs for universities. Following the independent review of the ARC Act in 2023 \citep[\textit{Trusting Australia's Ability},][]{sheil2023trusting} and the Australian Universities Accord \citep{accord2024final}, the Australian Government announced the discontinuation of the decennial ERA cycle in favor of developing a modern, data-driven national research assessment system that harnesses transparent global bibliometric data.

Similar transitions are underway internationally. The UK's Research Excellence Framework (REF) is continually reassessing the balance between expert peer assessment and algorithmic indicators \citep{wilsdon2015metric}, while international university league tables continue to proliferate. However, conventional bibliometric indicators remain fraught with systemic limitations when deployed for national evaluation:

\begin{enumerate}
    \item \textbf{Size and Volume Bias}: Raw publication counts and aggregate citation totals systematically advantage large, multi-faculty institutions over smaller, specialized, or regional universities \citep{docampo2014internal, docampo2015effect, docampo2017academic}. Size-dependent measures reward scale rather than research intensity or quality per output.
    \item \textbf{Journal Impact Factor (JIF) Fallacy}: Assigning value to an individual paper based exclusively on the journal in which it appears ignores the highly skewed distribution of citations within journals and overlooks the prestige and intellectual environment of the authoring institutions \citep{seglen1997why, hicks2015leiden}.
    \item \textbf{Field Incommensurability}: Citation densities, publication velocity, and co-authorship customs differ drastically between fields---for example, between clinical medicine and the humanities. Comparing raw citation counts across fields without rigorous normalization produces severe disciplinary distortions.
    \item \textbf{Gaming and Insularity}: Unipartite citation measures can be distorted by self-citation networks, citation clubs, and regional insularity, obscuring genuine scholarly influence.
\end{enumerate}

To overcome these structural defects, this paper presents an evidence-based framework for national research evaluation built upon a \textbf{bipartite spectral ranking model} \citep{docampo2025screening}. Instead of viewing journals or institutions in isolation, our approach formalizes scientific influence as a mutual reinforcement network between publishing channels (sources, $\mathcal{S}$) and research institutions ($\mathcal{I}$). High-prestige journals are those cited by leading research institutions, and high-prestige institutions are those that publish frequently in, and are cited by, influential journals. 

By applying a one-mode Breiger projection \citep{breiger1974duality} and solving for the dominant Perron eigenvector, we derive an intrinsic, size-independent influence intensity metric, $v$, scaled such that the global average across any discipline is exactly $1.0$. An institution with $v_I = 2.0$ in a given field produces research that attracts twice the global baseline prestige per unit of published work, irrespective of whether the university publishes 100 or 10,000 papers.

Furthermore, we ground this global model directly in national policy frameworks by creating an explicit concordance between the open scholarly graph provided by OpenAlex \citep{priem2022openalex} and the Australian and New Zealand Standard Research Classification \citep[ANZSRC 2020;][]{anzsrc2020} Fields of Research (FOR). We organize this evaluation at two complementary tiers:
\begin{itemize}
    \item A macro-level \textbf{5 Broad Research Areas (AREA5)}---separating computing and mathematical sciences as an independent analytical pillar from physical sciences and engineering, and properly aligning veterinary science with agricultural and animal sciences.
    \item A meso-level \textbf{26-discipline field tier}, mapping directly to 2-digit ANZSRC FOR divisions.
\end{itemize}

Applying this model to all 42 Australian Higher Education Providers (HEPs) across the 2020--2024 census window reveals distinct institutional profiles, demonstrating that research intensity flourishes across diverse mission groups---including the Group of Eight (Go8), the Australian Technology Network (ATN), the Innovative Research Universities (IRU), and Regional Universities Network (RUN). Finally, extensive sensitivity and geopolitical bloc analyses demonstrate that the spectral ranking is remarkably robust against shifts in citation boundaries and parameter variations.
